% p322_1; source PDF page 322; manually verified editable table.
\begin{tikzpicture}[x=1cm,y=-1cm,inner sep=0pt,outer sep=0pt]
\path[use as bounding box] (0,0) rectangle (16.000,15.390);
\draw[line width=.4pt] (0.000,0.000) rectangle (2.000,0.815);
\node[align=center,text width=1.740cm,font={\linespread{1}\fontsize{10}{13.50}\selectfont\bfseries}] at (1.000,0.407) {构成要件};
\draw[line width=.4pt] (2.000,0.000) rectangle (6.050,0.815);
\node[align=center,text width=3.790cm,font={\linespread{1}\fontsize{10}{13.50}\selectfont\bfseries}] at (4.025,0.407) {要点};
\draw[line width=.4pt] (6.050,0.000) rectangle (16.000,0.815);
\node[align=center,text width=9.690cm,font={\linespread{1}\fontsize{10}{13.50}\selectfont\bfseries}] at (11.025,0.407) {范例};
\draw[line width=.4pt] (0.000,0.815) rectangle (2.000,2.330);
\node[align=center,text width=1.740cm,font={\linespread{1}\fontsize{10}{13.50}\selectfont}] at (1.000,1.573) {行为主体};
\draw[line width=.4pt] (2.000,0.815) rectangle (6.050,2.330);
\node[align=left,text width=3.790cm,font={\linespread{1}\fontsize{10}{13.50}\selectfont}] at (4.025,1.573) {不限于驾驶者};
\draw[line width=.4pt] (6.050,0.815) rectangle (16.000,2.330);
\node[align=left,text width=9.690cm,font={\linespread{1}\fontsize{10}{13.50}\selectfont}] at (11.025,1.573) {甲看到乙醉酒，仍将车交给乙驾驶，乙醉驾不慎撞死人。乙构成交通肇事罪，甲也构成交通肇事罪};
\draw[line width=.4pt] (0.000,2.330) rectangle (2.000,5.545);
\node[align=center,text width=1.740cm,font={\linespread{1}\fontsize{10}{13.50}\selectfont}] at (1.000,3.938) {行为地点};
\draw[line width=.4pt] (2.000,2.330) rectangle (6.050,5.545);
\node[align=left,text width=3.790cm,font={\linespread{1}\fontsize{10}{13.50}\selectfont}] at (4.025,3.938) {公共交通领域};
\draw[line width=.4pt] (6.050,2.330) rectangle (16.000,5.545);
\node[align=left,text width=9.690cm,font={\linespread{1}\fontsize{10}{13.50}\selectfont}] at (11.025,3.938) {（1）在公共交通领域外，驾驶车辆过失致人死亡，构成过失致人死亡罪。\\{}(2）公共交通领域包括公路交通、水上交通。（2025年试题）\\{}（3）铁路事故，定铁路运营安全事故罪。\\{}（4）飞行事故，定重大飞行事故罪};
\draw[line width=.4pt] (0.000,5.545) rectangle (2.000,7.660);
\node[align=center,text width=1.740cm,font={\linespread{1}\fontsize{10}{13.50}\selectfont}] at (1.000,6.603) {肇事行为};
\draw[line width=.4pt] (2.000,5.545) rectangle (6.050,7.660);
\node[align=left,text width=3.790cm,font={\linespread{1}\fontsize{10}{13.50}\selectfont}] at (4.025,6.603) {肇事行为是过失行为};
\draw[line width=.4pt] (6.050,5.545) rectangle (16.000,7.660);
\node[align=left,text width=9.690cm,font={\linespread{1}\fontsize{10}{13.50}\selectfont}] at (11.025,6.603) {甲驾车想闯红灯，认为不会出事，不慎撞死行人。甲闯红灯是故意为之（行政法上的故意），但对死亡结果是过失心态（刑法上的过失)};
\draw[line width=.4pt] (0.000,7.660) rectangle (2.000,10.875);
\node[align=center,text width=1.740cm,font={\linespread{1}\fontsize{10}{13.50}\selectfont}] at (1.000,9.268) {实害结果};
\draw[line width=.4pt] (2.000,7.660) rectangle (6.050,10.875);
\node[align=left,text width=3.790cm,font={\linespread{1}\fontsize{10}{13.50}\selectfont}] at (4.025,9.268) {死亡1人；重伤3人；重伤1人，并有严重情节（酒驾、吸毒驾驶、无照驾驶、严重超载、肇事后逃逸)};
\draw[line width=.4pt] (6.050,7.660) rectangle (16.000,10.875);
\node[align=left,text width=9.690cm,font={\linespread{1}\fontsize{10}{13.50}\selectfont}] at (11.025,9.268) {这里的“肇事后逃逸”被作为本罪成立条件后，便不能再作为本罪的法定刑升格条件“肇事后逃逸，处3年至7年有期徒刑”，否则属于重复评价、重复处罚};
\draw[line width=.4pt] (0.000,10.875) rectangle (2.000,15.390);
\node[align=center,text width=1.740cm,font={\linespread{1}\fontsize{10}{13.50}\selectfont}] at (1.000,13.133) {因果关系};
\draw[line width=.4pt] (2.000,10.875) rectangle (6.050,15.390);
\node[align=left,text width=3.790cm,font={\linespread{1}\fontsize{10}{13.50}\selectfont}] at (4.025,13.133) {运用因果关系原理。特别是，实害结果必须是规范保护范围内的结果};
\draw[line width=.4pt] (6.050,10.875) rectangle (16.000,15.390);
\node[align=left,text width=9.690cm,font={\linespread{1}\fontsize{10}{13.50}\selectfont}] at (11.025,13.133) {甲驾驶没有年检的车辆，但该车并无故障。乙横穿高速公路，甲刹车不及撞死乙。“禁止驾驶未年检的车辆”，这个规范要防止的结果是，因车辆故障导致交通事故。现在乙死亡这个结果不是这个规范所能防止的结果，不能归属给甲。甲不构成交通肇事罪。\\{}[注意］不能将行政法上的违法直接等同于刑事违法，不能将行政法上的责任直接等同于刑事责任};
\end{tikzpicture}
